\documentclass[11pt,letterpaper]{article}
\usepackage[utf8]{inputenc}
\usepackage[spanish]{babel}
\usepackage[T1]{fontenc}
\usepackage{lmodern}
\usepackage{geometry}
\geometry{top=2cm, bottom=2.5cm, left=2cm, right=2cm}
\usepackage{listings}
\usepackage{xcolor}
\usepackage{amsmath}
\usepackage{amssymb}
\usepackage{booktabs}
\usepackage{fancyhdr}
\usepackage{titlesec}
\usepackage{graphicx}
\usepackage{hyperref}

% --- CONFIGURACIÓN DE COLORES ACADÉMICOS ---
\definecolor{darkblue}{rgb}{0.0, 0.2, 0.4}
\definecolor{codebg}{rgb}{0.97, 0.97, 0.99}
\definecolor{commentgreen}{rgb}{0.0, 0.5, 0.0}
\definecolor{keywordblue}{rgb}{0.0, 0.0, 0.8}
\definecolor{stringpurple}{rgb}{0.5, 0.0, 0.5}

% --- DISEÑO DE PÁGINA ---
\titleformat{\section}{\color{darkblue}\normalfont\Large\bfseries}{\thesection}{1em}{}
\titleformat{\subsection}{\color{darkblue}\normalfont\large\bfseries}{\thesubsection}{1em}{}
\titleformat{\subsubsection}{\color{darkblue}\normalfont\normalsize\bfseries}{\thesubsubsection}{1em}{}

\pagestyle{fancy}
\fancyhf{}
\lhead{\color{gray}FACYT - Universidad de Carabobo}
\rhead{\color{gray}Arquitectura del Computador - MIPS 32}
\cfoot{\thepage}

% --- CONFIGURACIÓN ESTRICTA DE LISTINGS PARA MIPS 32 ---
\lstdefinelanguage{MIPS}{
    morekeywords={add,addi,addu,addiu,sub,subu,and,andi,or,ori,xor,xori,nor,sll,srl,sra,lw,sw,lb,sb,lbu,lh,sh,lhu,lui,beq,bne,j,jal,jr,slt,slti,sltu,syscall,li,la,move,b,blez,bgez,blt,bgt,mul,div,mflo,mfhi},
    morecomment=[l]{\#},
    morestring=[b]",
}

\lstset{
    language=MIPS,
    backgroundcolor=\color{codebg},
    basicstyle=\ttfamily\small,
    breaklines=true,
    keywordstyle=\color{keywordblue}\bfseries,
    commentstyle=\color{commentgreen}\itshape,
    stringstyle=\color{stringpurple},
    numbers=left,
    numberstyle=\tiny\color{gray},
    numbersep=8pt,
    frame=single,
    frameround=tttt,
    framexleftmargin=4mm,
    rulecolor=\color{gray!30},
    keepspaces=true,
    showstringspaces=false,
    inputencoding=utf8,
    extendedchars=true,
    literate={á}{{\'a}}1 {é}{{\'e}}1 {í}{{\'i}}1 {ó}{{\'o}}1 {ú}{{\'u}}1 {ñ}{{\~n}}1 {ñ}{{\~n}}1 {Ñ}{{\~N}}1 {Í}{{\'{I}}}1 {Ó}{{\'{O}}}1
}

\begin{document}

% --- PORTADA ---
\begin{center}
    {\large \textbf{UNIVERSIDAD DE CARABOBO}} \\
    {\large \textbf{FACULTAD DE CIENCIA Y TECNOLOGÍA (FACYT)}} \\
    {\large \textbf{DEPARTAMENTO DE COMPUTACIÓN}} \\
    \vspace{0.8cm}
    {\huge \color{darkblue} \textbf{Guia de Referencia de MIPS 32}} \\
    \vspace{0.4cm}
    {\large \textbf{Materia: Arquitectura del Computador (4to Semestre)}} \\
    \vspace{0.2cm}
    \textit{Guía Exhaustiva de Código Basada en Patterson \& Hennessy (4ta Edición)}
    \vspace{0.6cm}
\end{center}

\hrule height 1.5pt
\vspace{0.6cm}

\section{Estructura Organizativa de Registros en MIPS 32}
MIPS utiliza una arquitectura Registro-Registro (RISC). Para codificar correctamente cualquier algoritmo, es mandatorio respetar el contrato de preservación de registros establecido por el hardware.

\begin{table}[h!]
\centering
\small
\begin{tabular}{cclc}
\toprule
\textbf{Nombre} & \textbf{Número} & \textbf{Uso en Arquitectura / Convención de Software} & \textbf{¿Se preserva en llamadas?} \\
\midrule
\texttt{\$zero} & 0 & Constante cableada al valor cero (0) & No modificable \\
\texttt{\$at}   & 1 & Reservado temporalmente para el Ensamblador & No \\
\texttt{\$v0-\$v1} & 2-3 & Valores de retorno de funciones evaluadas & No \\
\texttt{\$a0-\$a3} & 4-7 & Parámetros de argumentos de entrada & No \\
\texttt{\$t0-\$t7} & 8-15 & Registros temporales volátiles & No \\
\texttt{\$s0-\$s7} & 16-23 & Registros salvados (Contexto local a proteger) & \textbf{Sí} \\
\texttt{\$t8-\$t9} & 24-25 & Registros temporales volátiles adicionales & No \\
\texttt{\$k0-\$k1} & 26-27 & Reservados de uso exclusivo del Núcleo del S.O. & No \\
\texttt{\$gp}   & 28 & Puntero de área global de variables estáticas & Sí \\
\texttt{\$sp}   & 29 & Puntero de Pila (\textit{Stack Pointer}) & \textbf{Sí} \\
\texttt{\$fp}   & 30 & Puntero de Marco de Activación (\textit{Frame Pointer}) & Sí \\
\texttt{\$ra}   & 31 & Dirección de Retorno de la función actual (\textit{Return Address}) & \textbf{Sí} \\
\bottomrule
\end{tabular}
\caption{Mapa de Convenciones del Banco de Registros en MIPS.}
\end{table}

\section{Estructura Esquelética de un Programa en MARS}
Todo archivo ejecutable `.asm` debe dividir los recursos de memoria lógica mediante directivas claras del compilador.

\begin{lstlisting}[caption={Plantilla estructural limpia para desarrollo en laboratorios.}]
#----------------------------------------------------------------------------------
# AUTOR: Gabriel Vargas
# DESCRIPCIÓN: Estructura base para el desarrollo de proyectos MIPS
#----------------------------------------------------------------------------------

.data
    # Segmento de Memoria Estática (RAM)
    # Declaración de variables globales, cadenas ASCII y estructuras inicializadas
    mensaje_inicio: .asciiz "Iniciando ejecución del programa maestro...\n"
    salto_linea:    .asciiz "\n"
    alfa_numerico:  .byte 'M'
    constante_pi:   .word 314159   # Representación entera fija

.text
.globl main

main:
    # --- BLOQUE INICIAL DE CONTROL ---
    la $a0, mensaje_inicio
    li $v0, 4
    syscall                     # Impresión de bienvenida por consola
    
    # --- EJECUCIÓN ESPECÍFICA ---
    # (Pega aquí la lógica de tu algoritmo)

    # --- SALIDA REGLAMENTARIA Y LIMPIA ---
    li $v0, 10                  # Código del sistema para terminar el hilo actual
    syscall
\end{lstlisting}

\newpage

\section{Capítulo 2.2 a 2.5: Operaciones del Núcleo Aritmético y Lógico}
MIPS procesa tres tipos primarios de formatos de instrucción: Tipo R (Registros), Tipo I (Inmediatos) y Tipo J (Saltos). Aquí se muestran todas las variaciones posibles de manipulación aritmética y de bits.

\begin{lstlisting}[caption={Demostración total de operaciones de tipo aritmético-lógico.}]
# === 1. ARITMÉTICA DE ENTEROS ===
add  $t0, $s0, $s1       # Suma con desbordamiento controlado: $t0 = $s0 + $s1
addu $t1, $s0, $s1       # Suma sin signo (Ignora excepciones de overflow)
sub  $t2, $s2, $s3       # Resta básica: $t2 = $s2 - $s3
subu $t3, $s2, $s3       # Resta básica sin signo

# === 2. MANIPULACIÓN CON CONSTANTES (INMEDIATOS) ===
addi $t4, $s0, 256       # Suma inmediata con signo: $t4 = $s0 + 256
addiu $t5, $s0, -40      # Suma inmediata sin signo (Extiende el signo del inmediato)

# === 3. OPERACIONES DE BITS (MÁSCARAS) ===
and  $t6, $s1, $s2       # AND lógico bit a bit
andi $t7, $s1, 0x00FF    # AND inmediato: Extrae o aísla los 8 bits inferiores (Mascara de limpieza)
or   $s4, $t6, $t7       # OR lógico bit a bit
ori  $s5, $t6, 0x0F00    # OR inmediato: Enciende selectivamente grupos de bits
xor  $t8, $s1, $s2       # XOR (OR exclusivo): Detecta diferencias de bits
nor  $t9, $s0, $zero     # NOR contra el registro cero: Invierte todos los bits. Es el operador NOT.

# === 4. MULTIPLICACIÓN Y DIVISIÓN COMPLETA (Registros Especiales lo y hi) ===
mult $s0, $s1            # Multiplica $s0 * $s1. El resultado de 64 bits va a [hi, lo]
mflo $t0                 # Mueve la mitad inferior del producto desde 'lo' a $t0
mfhi $t1                 # Mueve la mitad superior del producto desde 'hi' a $t1

div  $s2, $s3            # Divide $s2 / $s3. Cociente se guarda en 'lo', Residuo en 'hi'
mflo $s4                 # $s4 = Cociente
mfhi $s5                 # $s5 = Residuo (Operación Módulo)

# === 5. DESPLAZAMIENTOS MÁQUINA (SHIFTS) ===
sll  $t0, $s1, 3         # Shift Left Logical: Desplaza 3 bits a la izquierda (Multiplica por 8)
srl  $t1, $s2, 2         # Shift Right Logical: Desplaza 2 bits a la derecha (Divide entre 4 sin signo)
sra  $t2, $s2, 2         # Shift Right Arithmetic: Desplaza conservando el bit de signo superior
\end{lstlisting}

\section{Capítulo 2.3 y 2.8: Modos de Direccionamiento de Memoria}
La memoria RAM lee y escribe usando bytes indexados de forma absoluta. Las palabras completas exigen alineación obligatoria de 4 bytes.

\begin{lstlisting}[caption={Lectura y almacenamiento de estructuras de datos discretas.}]
.data
    vector_enteros:   .word 105, 210, 315, 420
    vector_caracter:  .byte 'A', 'B', 'C', 'D', 'E'
    vector_cortos:    .half 12, 24, 36, 48

.text
main_memoria:
    la $s0, vector_enteros    # Cargar dirección base de la estructura de palabras (32 bits)
    la $s1, vector_caracter   # Cargar dirección base de la estructura de caracteres (8 bits)
    la $s2, vector_cortos     # Cargar dirección base de la estructura de cortos (16 bits)

    # LECTURAS (De memoria RAM a Registros CPU)
    lw  $t0, 0($s0)           # Carga vector_enteros[0] -> Desplazamiento base 0
    lw  $t1, 4($s0)           # Carga vector_enteros[1] -> Salta 4 bytes adelante
    lw  $t2, 12($s0)          # Carga vector_enteros[3] -> Salta 3 * 4 = 12 bytes adelante
    
    lbu $t3, 1($s1)           # Carga vector_caracter[1] como un byte sin signo (Carácter 'B')
    lh  $t4, 2($s2)           # Carga vector_cortos[1] -> Salta 1 * 2 = 2 bytes adelante

    # ESCRITURAS (De Registros CPU a memoria RAM)
    addi $t0, $t0, 100        # Modificar el valor local
    sw   $t0, 8($s0)          # Escribe el nuevo valor en vector_enteros[2] (Reemplaza el original)
    
    li   $t5, 'Z'
    sb   $t5, 0($s1)          # Sobrescribe vector_caracter[0] convirtiéndolo en 'Z'
\end{lstlisting}

\newpage

\section{Capítulo 2.6: Control de Flujo y Mapeo Algorítmico Base 1}

\subsection{Estructuras Condicionales Compuestas (If-Else / Selectores)}
Mapeo de la estructura lógica: \texttt{if (A < B) \{ C = A; \} else \{ C = B; \}} \\
Asignación de registros: \texttt{A = \$s0, B = \$s1, C = \$s2}.
\begin{lstlisting}[caption={Uso clásico de la comparación slt para bifurcaciones condicionales.}]
    slt $t0, $s0, $s1          # Si $s0 < $s1, entonces $t0 = 1; de lo contrario $t0 = 0
    beq $t0, $zero, ej_else    # Si $t0 == 0 (Es decir, A >= B), saltar al bloque Else
    
    # --- BLOQUE THEN ---
    move $s2, $s0              # C = A
    j fin_condicional          # Evitar la ejecución del bloque Else
    
ej_else:
    # --- BLOQUE ELSE ---
    move $s2, $s1              # C = B

fin_condicional:
\end{lstlisting}

\subsection{Estructuras de Bucles de Iteración (While / For)}
Código equivalente en C: \texttt{for (int i = 0; i < N; i++) \{ suma += vector[i]; \}} \\
Asignación: \texttt{i = \$t0, N = \$s0, suma = \$s1, dirección de vector = \$s2}.
\begin{lstlisting}[caption={Bucle indexado estándar con cálculo dinámico de offset.}]
    li $t0, 0                  # Inicializar el contador de ciclos i = 0
    li $s1, 0                  # Inicializar acumulador suma = 0

ciclo_for:
    slt $t1, $t0, $s0          # Verificar si i < N
    beq $t1, $zero, fin_for    # Si $t1 == 0 (i >= N), romper el bucle e ir al final
    
    # Cálculo del desplazamiento físico dentro del bucle
    sll $t2, $t0, 2            # $t2 = i * 4 (Cada palabra de entero ocupa 4 bytes)
    add $t3, $t2, $s2          # $t3 = Dirección específica física calculada (Base + Desplazamiento)
    
    lw  $t4, 0($t3)            # Cargar dinámicamente el elemento vector[i]
    add $s1, $s1, $t4          # suma += vector[i]
    
    addi $t0, $t0, 1           # Incrementar contador: i++
    j ciclo_for                # Evaluar la siguiente iteración

fin_for:
\end{lstlisting}

\subsection{Conversión Especial: Arreglos Base 1 (Algoritmos) a Base 0 (Hardware)}
En las evaluaciones teóricas de algoritmos en FACYT se suele asumir que los índices de los arreglos inician en \textbf{1}. En MIPS, las posiciones físicas de memoria siempre inician estrictamente en \textbf{0}.
El modelo matemático exacto para resolver este desajuste es:
\begin{equation}
\text{Direccion Fisica} = \text{Direccion Base} + [(\text{Indice Logico}_{\text{Base 1}} - 1) \times \text{TamaNo de la Palabra}]
\end{equation}

\begin{lstlisting}[caption={Rutina de adaptación segura para índices de pseudocódigo en Base 1.}]
# Parámetros: $s0 = Dirección Base del Arreglo en RAM, $s1 = Índice "i" en Base 1 deseado.
# Requerimiento: Leer el dato simulado arreglo[i]

addi $t0, $s1, -1          # Ajuste matemático inmediato: Convertir a Base 0 ($t0 = i - 1)
bltz $t0, error_limite     # Protección: Si el índice ajustado es menor a 0, abortar por violación

sll  $t1, $t0, 2           # Multiplicar el índice ajustado por 4 bytes (Tamaño de palabra)
add  $t2, $t1, $s0         # Calcular la dirección real combinada
lw   $s2, 0($t2)           # Cargar con éxito el dato equivalente al algoritmo teórico

j fin_acceso

error_limite:
    # Manejo preventivo de excepciones de desbordamiento de arreglos
fin_acceso:
\end{lstlisting}

\newpage

\section{Capítulo 2.7: Administración de Procedimientos y Gestión de la Pila}

\subsection{Estructura de Funciones Anidadas No-Hoja}
Una función No-Hoja ejecuta sub-llamadas internas mediante la instrucción `jal`. Esto causa la pérdida destructiva inmediata del contenido del registro `$ra`. Por lo tanto, la función debe crear su propio marco de pila seguro en memoria dinámica.

\begin{lstlisting}[caption={Plantilla estándar obligatoria para el prólogo y epílogo del manejo de la Pila.}]
#----------------------------------------------------------------------------------
# FUNCIÓN EJEMPLO: Procesa datos y realiza una llamada anidada interna
#----------------------------------------------------------------------------------
funcion_compleja:
    # === 1. PRÓLOGO DE LA FUNCIÓN ===
    addi $sp, $sp, -8          # Reservamos espacio para 2 palabras continuas (8 bytes)
    sw   $ra, 4($sp)           # Guardamos la dirección de retorno del llamador original
    sw   $s0, 0($sp)           # Salvamos el valor previo del registro $s0

    # === 2. CUERPO OPERATIVO ===
    move $s0, $a0              # Protegemos el primer argumento de entrada en $s0
    add  $a0, $s0, $a1         # Preparamos un parámetro modificado para la sub-función
    
    jal  sub_rutina_interna    # Ejecutar la llamada anidada. Al hacer esto se altera $ra de forma local.
                               # El resultado parcial de la sub-rutina se recibe en el registro $v0
                               
    mul  $v0, $v0, $s0         # Multiplicar el resultado parcial por el valor inicial guardado en $s0

    # === 3. EPÍLOGO DE LA FUNCIÓN ===
    lw   $s0, 0($sp)           # Restauramos el valor original e intacto de $s0 del llamador
    lw   $ra, 4($sp)           # Recuperamos la dirección de retorno REAL hacia main
    addi $sp, $sp, 8           # Liberamos los 8 bytes sumando de vuelta al puntero de pila
    
    jr   $ra                   # Retornar de forma exitosa y segura
\end{lstlisting}

\subsection{Implementación de un Algoritmo Recursivo Puro (Factorial)}
Demostración avanzada de recursión matemática controlada por pila.
\begin{lstlisting}[caption={Código máquina para el cálculo recursivo de N!.}]
# Argumento de Entrada: $a0 = Número entero N
# Retorno: $v0 = Resultado del factorial (N!)
factorial:
    addi $sp, $sp, -8          # Reservar espacio en la pila para el sub-árbol
    sw   $ra, 4($sp)           # Salvar dirección de retorno
    sw   $a0, 0($sp)           # Salvar el valor del argumento actual N
    
    slti $t0, $a0, 1           # Evaluar caso base: Es N < 1?
    beq  $t0, $zero, ir_recursio # Si N >= 1, continuar con las llamadas descendientes
    
    # --- CASO BASE ALCANZADO (0! = 1) ---
    li   $v0, 1                 # Cargar el resultado neutro del inicio
    addi $sp, $sp, 8           # Liberar pila del nivel actual sin leer memoria
    jr   $ra                   # Iniciar el retorno hacia arriba
    
ir_recursio:
    addi $a0, $a0, -1          # Decrementar el argumento de control: (N - 1)
    jal  factorial             # Llamada recursiva interna: factorial(N-1)
                               # Al regresar, $v0 contiene el resultado acumulado acumulado
                               
    lw   $a0, 0($sp)           # Restaurar el valor original de N asignado a este nivel
    lw   $ra, 4($sp)           # Restaurar la dirección de retorno de este nivel
    addi $sp, $sp, 8           # Liberar el espacio del marco de pila actual
    
    mul  $v0, $a0, $v0         # Multiplicación combinada: N * factorial(N - 1)
    jr   $ra                   # Volver al llamador del nivel superior
\end{lstlisting}

\newpage

\section{Algorítmica Avanzada para Evaluaciones y Proyectos}

\subsection{Ordenamiento de Arreglos (Procedimiento Completo Sort y Swap)}
El algoritmo de ordenamiento por intercambio (\textit{Bubble / Exchange Sort}) representa el ejemplo cumbre de un procedimiento No-Hoja con dos bucles condicionales aninados según el texto guía de Patterson \& Hennessy.

\begin{lstlisting}[caption={Código máquina MIPS completo para la rutina sort y sub-rutina swap.}]
#----------------------------------------------------------------------------------
# PROCEDIMIENTO PRINCIPAL DE ORDENAMIENTO: sort (Función No-Hoja)
# Argumentos: $a0 = Dirección Base del Arreglo, $a1 = Tamaño N
# Mapeo: $s0 = i, $s1 = j, $s2 = v (dirección base), $s3 = n (tamaño)
#----------------------------------------------------------------------------------
sort:
    # --- PRÓLOGO: Reservar espacio en la Pila (Stack Frame) ---
    addi $sp, $sp, -20         # Espacio para resguardar 5 palabras (20 bytes)
    sw   $ra, 16($sp)          # ¡VITAL! Guardar retorno a main
    sw   $s3, 12($sp)          # Guardar contexto de $s3
    sw   $s2, 8($sp)           # Guardar contexto de $s2
    sw   $s1, 4($sp)           # Guardar contexto de $s1
    sw   $s0, 0($sp)           # Guardar contexto de $s0

    # --- INICIALIZACIÓN ---
    move $s2, $a0              # Copiar dirección base de v para protegerla
    move $s3, $a1              # Copiar el tamaño N para protegerlo
    move $s0, $zero            # i = 0 (Inicializa el lazo externo)

# === PRIMER LAZO FOR (Externo) ===
for_externo_test:
    slt  $t0, $s0, $s3         # i < n?
    beq  $t0, $zero, exit_externo # Si $t0 == 0 (i >= n), ir al epílogo final
    
    addi $s1, $s0, -1          # Inicialización del lazo interno: j = i - 1

# === SEGUNDO LAZO FOR (Interno) ===
for_interno_test:
    # Condición Parte 1: j >= 0?
    slti $t0, $s1, 0           # Si j < 0 entonces $t0 = 1
    bne  $t0, $zero, exit_interno # LÓGICA INVERTIDA: Si j < 0, sale del lazo interno
    
    # Condición Parte 2: v[j] > v[j + 1]?
    sll  $t1, $s1, 2           # $t1 = j * 4 (Convertir índice a bytes)
    add  $t2, $s2, $t1         # $t2 = v + (j * 4) -> Dirección física de v[j]
    
    lw   $t3, 0($t2)           # $t3 = v[j]
    lw   $t4, 4($t2)           # $t4 = v[j + 1] (4 bytes más adelante)
    
    slt  $t0, $t4, $t3         # v[j + 1] < v[j]? (Equivale a v[j] > v[j+1])
    beq  $t0, $zero, exit_interno # LÓGICA INVERTIDA: Si no es mayor, rompe el ciclo

    # --- CUERPO DEL FOR INTERNO (Llamada a Swap) ---
    move $a0, $s2              # Parámetro 1: Dirección base v
    move $a1, $s1              # Parámetro 2: Índice j actual
    jal  swap                  # Saltar a sub-rutina swap. Al regresar ejecuta la sig. línea
    
    # --- ACTUALIZACIÓN DEL FOR INTERNO ---
    addi $s1, $s1, -1          # j--
    j    for_interno_test      # Volver a evaluar el lazo interno

exit_interno:
    # --- ACTUALIZACIÓN DEL FOR EXTERNO ---
    addi $s0, $s0, 1           # i++
    j    for_externo_test      # Volver a evaluar el lazo externo

exit_externo:
    # --- EPÍLOGO: Restaurar la Pila y Retornar ---
    lw   $s0, 0($sp)           # Restaurar $s0 original del llamador
    lw   $s1, 4($sp)           # Restaurar $s1
    lw   $s2, 8($sp)           # Restaurar $s2
    lw   $s3, 12($sp)          # Restaurar $s3
    lw   $ra, 16($sp)          # Recuperar el verdadero $ra que apunta al main
    addi $sp, $sp, 20          # Liberar de forma exacta los 20 bytes
    jr   $ra                   # Retornar de forma segura


#----------------------------------------------------------------------------------
# PROCEDIMIENTO AUXILIAR: swap (Función Hoja)
# Argumentos: $a0 = Dirección Base v, $a1 = Índice j
#----------------------------------------------------------------------------------
swap:
    sll  $t0, $a1, 2           # $t0 = j * 4
    add  $t1, $a0, $t0         # $t1 = Dirección física de v[j]
    
    lw   $t2, 0($t1)           # $t2 = v[j]
    lw   $t3, 4($t1)           # $t3 = v[j + 1]
    
    sw   $t3, 0($t1)           # v[j] = v[j + 1] (Escritura cruzada)
    sw   $t2, 4($t1)           # v[j + 1] = v[j]
    
    jr   $ra                   # Retorna directo a la línea interna de sort
\end{lstlisting}

\subsection{Acceso y Mapeo de Matrices Bidimensionales (Filas y Columnas)}
Para procesar una matriz estática guardada en formato lineal por filas (\textit{Row-Major Order}), la fórmula matemática para calcular el desplazamiento del elemento \texttt{Matriz[fila][columna]} es:
\begin{equation}
\text{Offset en Bytes} = (\text{Fila} \times \text{Numero de Columnas Totales} + \text{Columna}) \times \text{TamaNo del Dato}
\end{equation}

\begin{lstlisting}[caption={Rutina de extracción de elementos en matrices dinámicas de enteros.}]
.data
    # Matriz simulada de 3x4 (3 filas, 4 columnas) con enteros secuenciales
    matriz_datos: .word  11, 12, 13, 14    # Fila 0
                  .word  21, 22, 23, 24    # Fila 1
                  .word  31, 32, 33, 34    # Fila 2
    columnas_tot: .word  4                 # Constante de ancho estructural

.text
main_matrices:
    la  $s0, matriz_datos      # Dirección base de la estructura en RAM
    lw  $s1, columnas_tot      # $s1 = 4 (Ancho total de la matriz)
    
    # Coordenadas lógicas del elemento que deseamos extraer (Ej: Fila 1, Columna 3 -> Valor 24)
    li  $s2, 1                 # Fila objetivo = 1
    li  $s3, 3                 # Columna objetivo = 3 (Buscamos matriz_datos[1][3] -> 24)
    
    # Implementación de la fórmula física de direccionamiento
    mul $t0, $s2, $s1          # $t0 = Fila objetivo * Columnas totales (1 * 4 = 4)
    add $t0, $t0, $s3          # $t0 = Linealización de coordenadas ($t0 = 4 + 3 = 7)
    
    sll $t1, $t0, 2            # Convertir a bytes multiplicando por 4 ($t1 = 7 * 4 = 28 bytes de offset)
    add $t2, $t1, $s0          # $t2 = Dirección física exacta en la RAM
    
    lw  $s4, 0($t2)            # $s4 = Contenido extraído de la celda deseada (Valor = 24)
\end{lstlisting}

\section{Catálogo Completo de Servicios de Interrupción (Syscalls)}
Carga el identificador en \texttt{\$v0} y pasa los parámetros requeridos antes de activar el servicio.

\begin{table}[h!]
\centering
\small
\begin{tabular}{ccll}
\toprule
\textbf{Servicio} & \textbf{Código (\$v0)} & \textbf{Argumentos de Entrada requeridos} & \textbf{Resultado del Sistema} \\
\midrule
Imprimir Entero & 1 & \texttt{\$a0} = Contiene el valor numérico entero & Muestra el número en consola \\
Imprimir Cadena & 4 & \texttt{\$a0} = Dirección base de la cadena (\texttt{.asciiz}) & Muestra el texto en consola \\
Leer Entero & 5 & Ninguno & Captura entrada; resultado en \texttt{\$v0} \\
Leer Cadena & 8 & \texttt{\$a0} = Dirección del Buffer, \texttt{\$a1} = Tamaño límite & Almacena el texto en memoria \\
Asignar Memoria & 9 & \texttt{\$a0} = Cantidad de bytes requeridos & Devuelve dirección asignada en \texttt{\$v0} \\
Terminar Programa & 10 & Ninguno & Cierra el hilo del programa de forma limpia \\
\bottomrule
\end{tabular}
\caption{Servicios de llamadas de consola indispensables para laboratorios en MARS.}
\end{table}

\end{document}